\section{Refinamento diádico adaptativo certificado}
\label{sec:adaptive}

Três amostras de uma função não fornecem, sem hipótese adicional, um limitante do erro em todo o intervalo. A garantia adaptativa será baseada em majorantes locais certificados de $|F''|$, não em diferenças segundas amostrais.

\subsection{Fechamento operacional por aritmética intervalar}

O majorante local pode ser construído, em vez de postulado, para uma gramática
explícita de funções. Seja $\mathcal E$ a menor classe de expressões que contém
$x$ e as constantes racionais e é fechada por soma, subtração, produto,
recíproco, $\exp$, $\sin$, $\cos$ e composição. Um recíproco só é admissível
em intervalos nos quais seu denominador não se anula. Essa classe contém os
polinômios, as funções racionais sem polos no domínio e as funções elementares
usadas nos experimentos.

Aqui composição significa substituição sintática: em $u\circ v$, cada folha
$x$ da árvore de $u$ é substituída pela árvore de $v$. Portanto toda expressão
composta volta a ser uma árvore das operações-base acima; composição não cria
um operador adicional nem um caso omitido na indução estrutural.

Escreva $\mathbb I_{\mathbb Q}$ para os intervalos fechados com extremos
racionais. A implementação pública
\begin{equation}
 \operatorname{BoundSecondDerivative}(E,J;P,d)
   =[\underline m_J,\overline m_J]\in\mathbb I_{\mathbb Q}
 \label{eq:bound-second-interface}
\end{equation}
usa precisão transcendental $P$ e profundidade de subdivisão $d$. Ela propaga
simultaneamente intervalos para $(E,E',E'')$, isto é, 2-jatos intervalares. Por
exemplo,
\begin{align*}
 (uv)''&=u''v+2u'v'+uv'',\\
 (u^{-1})''&=2(u')^2u^{-3}-u''u^{-2},\\
 (e^u)''&=e^u\bigl((u')^2+u''\bigr),\\
 (\sin u)''&=-\sin(u)(u')^2+\cos(u)u'',\\
 (\cos u)''&=-\cos(u)(u')^2-\sin(u)u''.
\end{align*}
As operações racionais são exatas. Para $\exp$, a redução do argumento a
$|z|\le1/2$ é seguida por uma soma de Taylor cujo resto é majorado por uma
série geométrica racional; para $\sin$ e $\cos$, usa-se o resto de Lagrange
racional e a interseção com $[-1,1]$. Depois, elevam-se $P$ e $d$ quando se
deseja reduzir a sobrestimação. Essa construção segue o princípio de inclusão
da aritmética intervalar validada \citep{moore2009,silva2025}.

\begin{lemma}[restos transcendentais com extremos racionais]
\label{lem:rational-transcendental-enclosures}
Se $I=[a,b]\in\mathbb I_{\mathbb Q}$, as rotinas descritas acima produzem
intervalos de extremos racionais que contêm $\exp(I)$, $\sin(I)$ e $\cos(I)$.
Mais explicitamente, para $0\le z\le1/2$ e $n\ge0$,
\begin{equation}
 0\le e^z-\sum_{j=0}^{n}\frac{z^j}{j!}
 \le
 \frac{z^{n+1}}{(n+1)!}
 \frac{1}{1-z/(n+2)}.
 \label{eq:exp-rational-tail}
\end{equation}
Se $R=\max\{|a|,|b|\}$ e
\begin{align*}
 P_n^{\sin}(z)&=\sum_{j=0}^{n}(-1)^j\frac{z^{2j+1}}{(2j+1)!},&
 P_n^{\cos}(z)&=\sum_{j=0}^{n}(-1)^j\frac{z^{2j}}{(2j)!},
\end{align*}
então, para todo $z\in I$,
\begin{align}
 |\sin z-P_n^{\sin}(z)|
 &\le \frac{R^{2n+2}}{(2n+2)!},
 &
 |\cos z-P_n^{\cos}(z)|
 &\le \frac{R^{2n+1}}{(2n+1)!}.
 \label{eq:trig-rational-remainders}
\end{align}
Avaliar os polinômios em $I$, inflar pelos respectivos restos e intersectar
com $[-1,1]$ preserva a inclusão.
\end{lemma}

\begin{proof}
Na série de $e^z$, o primeiro termo omitido é
$t_{n+1}=z^{n+1}/(n+1)!$. Para $j\ge n+1$, a razão entre termos consecutivos
é $z/(j+1)\le z/(n+2)<1$; a cauda é, portanto, majorada pela série geométrica
de primeiro termo $t_{n+1}$ e razão $z/(n+2)$, o que dá
\eqref{eq:exp-rational-tail}. Para um argumento racional $x$ arbitrário,
escolha $r$ com $z=|x|/2^r\le1/2$, aplique a inclusão anterior e eleve o
intervalo ao quadrado $r$ vezes; se $x<0$, tome o recíproco do intervalo
positivo obtido para $-x$. Em $I=[a,b]$, a monotonicidade de $e^x$ permite
usar o extremo inferior da inclusão em $a$ e o superior da inclusão em $b$.

As duas desigualdades em \eqref{eq:trig-rational-remainders} são o resto de
Lagrange dos polinômios de Maclaurin de graus $2n+1$ e $2n$, pois todas as
derivadas de seno e cosseno têm módulo no máximo $1$. A avaliação intervalar
contém a imagem de cada polinômio; inflar pelo raio do resto contém a função.
Como seno e cosseno já pertencem a $[-1,1]$, a interseção não remove nenhum
valor possível. Todas as operações, fatoriais e extremos envolvidos são
racionais.
\end{proof}

\begin{proposition}[correção do majorante intervalar]
\label{prop:interval-majorant}
Se $E\in\mathcal E$ define uma função $C^2$ em $J$ e a rotina
\eqref{eq:bound-second-interface} termina com sucesso, então
\begin{equation}
 E''(x)\in[\underline m_J,\overline m_J]\quad(x\in J),
 \qquad
 M_{\rm int}(J):=\max\{|\underline m_J|,|\overline m_J|\}
 \ge \|E''\|_{\infty,J}.
 \label{eq:interval-majorant}
\end{equation}
Se algum intervalo de denominador contém zero mesmo após a subdivisão
permitida, a rotina devolve falha e nenhum número é aceito como certificado.
Em particular, sua garantia nunca depende de amostragem.
\end{proposition}

\begin{proof}
Para constantes e para $x$, a inclusão é imediata. Soma, produto e recíproco
preservam inclusão pelas operações elementares de intervalos; o recíproco só
é executado quando zero está fora do intervalo do denominador. As fórmulas de
2-jatos acima e o \cref{lem:rational-transcendental-enclosures}
dão o passo de indução estrutural para toda expressão de $\mathcal E$.
Pela normalização sintática declarada antes do lema, uma composição genérica
já foi expandida nessa árvore e não acrescenta um caso independente.
Subdividir $J$ preserva a cobertura, e tomar o fecho convexo dos intervalos
obtidos em suas partes contém a união das imagens de $E''$. Isso prova a
primeira inclusão e, tomando módulos, \eqref{eq:interval-majorant}. A regra de
falha impede que uma divisão não validada produza uma falsa garantia.
\end{proof}

A certificação cobre as expressões para as quais os orçamentos $(P,d)$ fecham
a propagação intervalar sem polos. Formas naturais podem ser largas pelo
problema de dependência; subdivisão, formas centradas e Taylor intervalar
fornecem mecanismos explícitos de aperto. A implementação cobre a gramática
acima por subdivisão e Taylor racional e registra a largura, o número de caixas
e a profundidade usados.

\begin{algorithm}[refinamento certificado]\label{alg:adaptive}
Dados $f\in C^2[\alpha,\beta]$, tolerâncias $\varepsilon_I,\varepsilon_R>0$ e a extensão $F$ de \eqref{eq:c2-extension}, use, para toda célula $J$, um majorante validado
\begin{equation}
 \|F''\|_{\infty,J}\le M(J)\le M_*<\infty.
 \label{eq:majorant-assumption}
\end{equation}
Quando $F$ é fornecida por uma expressão de $\mathcal E$, tome
$M(J)=M_{\rm int}(J)$ de \eqref{eq:interval-majorant}; se a certificação
falhar, aborte sem emitir uma montagem garantida. Para uma função $C^2$
arbitrária fora dessa interface, \eqref{eq:majorant-assumption} continua sendo
um contrato de entrada explícito.
\begin{enumerate}[label=(\roman*),leftmargin=2.1em]
 \item Inicie com uma cobertura diádica finita de $[\alpha,\beta]$.
 \item Para cada folha $J$, calcule
 \begin{equation}
  \eta_{\mathrm{cert}}(J)=\frac{|J|^2M(J)}8.
  \label{eq:certificate}
 \end{equation}
 \item Enquanto existir folha com $\eta_{\mathrm{cert}}(J)>\varepsilon_I$, substitua-a por seus dois filhos.
 \item Nos nós finais, escolha racionais compartilhados $v_t$ com $|v_t-F(t)|<\varepsilon_R$ e valores adjacentes distintos.
 \item Emita, em cada folha, o endereço do átomo que interpola seus dois valores nodais.
\end{enumerate}
\end{algorithm}

\begin{theorem}[correção, término e tamanho]\label{thm:adaptive}
O \cref{alg:adaptive} termina e produz $g\in\Aclass([\alpha,\beta])$ com
\begin{equation}
 \|f-g\|_\infty\le\varepsilon_I+\varepsilon_R.
 \label{eq:adaptive-error}
\end{equation}
Se a cobertura inicial tem $R_0$ raízes de largura máxima $h_0$ e $M_*>0$, nenhuma ramificação ultrapassa
\begin{equation}
 r_*=\max\left\{0,
 \left\lceil\frac12\log_2\!\left(\frac{h_0^2M_*}{8\varepsilon_I}\right)\right\rceil
 \right\}
 \label{eq:depth}
\end{equation}
níveis, e há no máximo $R_0\,2^{r_*}$ folhas. Se $M_*=0$, pode-se tomar $r_*=0$.
\end{theorem}

\begin{proof}
Após $r$ subdivisões, uma folha descendente de uma raiz tem $|J|\le h_0\,2^{-r}$ e, por \eqref{eq:majorant-assumption},
\[
 \eta_{\mathrm{cert}}(J)\le\frac{h_0^2M_*}{8}4^{-r}.
\]
O lado direito é no máximo $\varepsilon_I$ para $r\ge r_*$; logo nenhuma folha é subdividida além dessa profundidade. A árvore é finita e cada raiz gera no máximo $2^{r_*}$ folhas.

Na partição final, \eqref{eq:interpolation-remainder} e \eqref{eq:majorant-assumption} implicam
\[
 \|F-P_\Lambda F\|_{\infty,J}
 \le\frac{|J|^2}{8}\|F''\|_{\infty,J}
 \le\eta_{\mathrm{cert}}(J)\le\varepsilon_I.
\]
A racionalização nodal produz, pelos \cref{thm:dictionary,thm:injective}, uma montagem de átomos de endereço único e com $\|P_\Lambda F-g\|_\infty<\varepsilon_R$. Como $F=f$ no intervalo original, a desigualdade triangular prova \eqref{eq:adaptive-error}.
\end{proof}

O \cref{thm:adaptive} fornece a garantia de término e de tamanho usada na
comparação posterior com as constantes diádicas; resultados recentes sobre
ordens assintóticas exatas em aproximação adaptativa oferecem um contexto mais
amplo para esse tipo de contagem \citep{kessebohmerniemann2025}.

\begin{proposition}[calibração exata para curvatura constante]
\label{prop:adaptive-dyadic-oscillation}
Suponha que $[\alpha,\beta]=J_0$ seja uma única célula diádica inicial,
$F''\equiv\kappa\ne0$ em $J_0$ e $M(J)=|\kappa|$ em todo descendente.
Pondo
\[
 A=\frac{|J_0|^2|\kappa|}{8},
\]
o número de folhas produzido com tolerância $\varepsilon_I$ é
\begin{equation}
 N(\varepsilon_I)=2^{r(\varepsilon_I)},\qquad
 r(\varepsilon_I)=
 \max\!\left\{0,\left\lceil\log_4\!\left(\frac{A}{\varepsilon_I}\right)\right\rceil\right\}.
 \label{eq:constant-curvature-leaves}
\end{equation}
Consequentemente,
\begin{equation}
 \liminf_{\varepsilon_I\downarrow0}\sqrt{\varepsilon_I}\,N(\varepsilon_I)
 =\frac1{\sqrt8}\int_\alpha^\beta |f''(x)|^{1/2}\dd x,
 \qquad
 \limsup_{\varepsilon_I\downarrow0}\sqrt{\varepsilon_I}\,N(\varepsilon_I)
 =\frac1{\sqrt2}\int_\alpha^\beta |f''(x)|^{1/2}\dd x.
 \label{eq:constant-curvature-oscillation}
\end{equation}
Em particular, a quantização diádica pode introduzir um fator dois e não se
pode atribuir ao algoritmo, em geral, a constante assintótica dos nós livres.
\end{proposition}

\begin{proof}
Após $r$ subdivisões uniformes, cada folha tem largura $|J_0|2^{-r}$ e
certificado $A4^{-r}$. Como todos os certificados de um mesmo nível são
iguais, o algoritmo para no primeiro $r$ com $A4^{-r}\le\varepsilon_I$, o
que prova \eqref{eq:constant-curvature-leaves}. Para
$A4^{-r}\le\varepsilon_I<A4^{-(r-1)}$, tem-se
\[
 \sqrt{A}\le \sqrt{\varepsilon_I}\,2^r<2\sqrt{A}.
\]
Os extremos inferior e superior são atingidos, respectivamente, nos limiares
$\varepsilon_I=A4^{-r}$ e por aproximação pela esquerda aos limiares
$A4^{-(r-1)}$. Por fim,
$\int_\alpha^\beta |f''|^{1/2}=|J_0|\sqrt{|\kappa|}$.
\end{proof}

\begin{remark}[diferença segunda não é certificado]\label{rem:second-difference}
Se $x_0=(x_L+x_R)/2$, então
\[
 |f(x_0)-P_Jf(x_0)|=\frac12|f(x_L)-2f(x_0)+f(x_R)|.
\]
Isso mede exatamente o erro no centro, mas não o máximo no intervalo. Para $f(x)=x^3$ em $[-1,1]$, $P_Jf(x)=x$ e a diferença segunda central é zero, enquanto
\[
 \|x^3-x\|_\infty=\frac{2}{3\sqrt3}\approx0{,}3849.
\]
Portanto, o teste de três pontos é diagnóstico, não garantia.
\end{remark}

O uso de majorantes rigorosos distingue o \cref{alg:adaptive} de heurísticas amostrais. Essa exigência é compatível com a literatura de aproximação adaptativa em árvores, na qual a confiabilidade depende de indicadores ou estimadores controláveis \citep{cohen2001,binev2004}.
